\chapter{Ethical Issues}

This project investigates an emotional AI engine that operates on the basis of latent representations of emotions. It will utilise recent concept-based reasoning paradigms (e.g.\ latent and implicit inference models) and use human judgement to evaluate the output. While the project does not involve invasive experiments or high-risk applications, it raises several ethical and legal considerations.

\section{Ethical Considerations}
The main ethical considerations relate to the involvement of human participants in the evaluation phase. Human evaluators are expected to assess the outputs of the model rather than disclose personal experiences or emotional states. Participation should be voluntary, and they should be free to withdraw at any time without consequence. The data collected shall be kept to a minimum (e.g.\ numerical evaluations and brief qualitative judgments only) and no personally identifiable information is requested.

The second ethical consideration concerns the nature of the system under study. This project explores the expression and regulation of emotions in AI systems. There is a potential ethical risk that these technologies, if misused, could be perceived as enabling emotional manipulation and deceptive emotional signals. However, the manipulation techniques addressed in this study are essentially identical to the widely adopted usage patterns of LLMs today. For example, instructions such as ``Please write more carefully'' or ``Could you rephrase this in a gentler manner?'' are already routinely provided as prompts, and the models adjust the emotional tone of their output accordingly. In other words, the significance of this project lies not in endowing models with new capabilities for emotional manipulation, but in visualising the structures of emotional expression already present within the models. The target of manipulation is not human emotion but the latent representations within the model. The value of this research lies in treating the ambiguous tone adjustments already occurring as a subject of direct study, rather than leaving them as a black box.

\section{Legal Considerations}

The project will also use open source software, pre-trained models (e.g.\ GPT-2) and public datasets. These resources will be used strictly in accordance with their respective licences and for non-commercial academic research purposes only. No proprietary, restricted or export-controlled software or information will be generated.